\section{Related Work}
\label{sec:2}

This study brings together four main research areas: dependability, performance, and sustainability in distributed software systems; static code and system analysis; software quality analysis and multi-criteria evaluation; and graph representation learning with explainable AI. Together, these fields provide the basis for predicting cascading vulnerabilities from Architecture-as-Code descriptors before setting up any infrastructure.

\subsection{Dependability, Performance, and Sustainability in Distributed Software Systems}
\label{sec:2.1}

Publish-subscribe (pub-sub) and asynchronous event-driven models separate communicating entities by space, time, and synchronization, enabling scalable, high-throughput systems \cite{eugster2003faces}. Modern middleware standards like ROS~2 \cite{macenski2022robot}, Apache Kafka \cite{kreps2011kafka}, DDS \cite{omg2015dds}, and MQTT \cite{oasis2019mqtt} manage these exchanges using detailed Quality-of-Service (QoS) policies. In cloud-native microservice meshes and distributed AI or LLM serving systems, asynchronous message passing and queueing are the main ways components communicate, affecting latency, throughput, and hardware use.

Previous research on dependability and performance has focused on runtime methods such as dynamic consensus, broker clustering, backpressure throttling, autoscaling, and automated failover. It has also explored chaos engineering and runtime fault injection \cite{basiri2016chaos}, which intentionally disrupt live clusters to study degradation and recovery. While runtime testing is realistic, it requires a working cluster, can disrupt services, and uses significant computing resources. These factors make it less practical during early design or commit-level CI/CD, and it is one of the development-time computations examined in green software engineering \cite{calero2015green,lannelongue2021green,verdecchia2023greenai,schmidt2021codecarbon,noureddine2021joularjx}. By contrast, this study examines dependability using Architecture-as-Code descriptors before setting up any infrastructure.

A long history of research has predicted dependability from architectural descriptions. For example, Cheung’s absorbing-Markov-chain model \cite{cheung1980user} calculates system reliability using component reliabilities and control flow graphs. Goseva-Popstojanova and Trivedi \cite{gosevapopstojanova2001architecture} organized later state-based, path-based, and additive models, as reviewed by Immonen and Niemel{\"a} \cite{immonen2008survey}. Model-driven frameworks like the Palladio Component Model \cite{becker2009palladio} and layered queueing networks \cite{franks2009layered} use component specifications to predict performance and reliability. Annotation-based methods, such as the AADL Error Model Annex \cite{delange2014aadl}, create fault trees from declared error states. However, these approaches cover broader topics than this study and need operational profiles and failure rates that are not available at commit time. Instead, SaG focuses on a specific question using only manifests: which components’ failures spread the furthest in the declared topology?

Many studies find faults in microservices by using operational telemetry. For example, Seer \cite{gan2019seer} and Sage \cite{gan2021sage} predict QoS problems from hardware counters and traces. MicroRCA \cite{wu2020microrca} and TraceRCA \cite{li2021tracerca} find root causes using service-dependency graphs. DeepTraLog \cite{zhang2022deeptralog}, Eadro \cite{lee2023eadro}, and MicroCause \cite{meng2020microcause} use graph neural networks on traces, logs, and metrics, as Zhang et al.~\cite{zhang2024failurediag} review. Zhou et al.~\cite{zhou2021fault} found that cascading outages in synchronous microservices often stem from thread-pool exhaustion, RPC timeouts, and repeated retry storms along call trees. In pub-sub systems, failures spread through broker queue saturation and message starvation. Because these methods require a running cluster to produce runtime telemetry, SaG works before deployment by analyzing static manifests before any code runs.

The question SaG asks --- which components' failure reaches furthest --- is also the question of software change impact analysis \cite{bohner1996impact}, which traces how a modification propagates through dependency structure, and of dependency management in microservice fleets \cite{esparrachiari2018tracking}, where cyclic and undeclared dependencies are a known outage source. Ranking components by predicted impact, rather than classifying them, relates to learning-to-rank formulations of defect prediction \cite{yang2015ltr}, which optimize the ranking measure directly as our listwise loss does (Section~\ref{sec:4.2}).

\subsection{Static Code Analysis (SCA) vs. Static System Analysis (SSA)}
\label{sec:2.2}

Traditional \textbf{Static Code Analysis (SCA)} tools (e.g., SonarQube \cite{sonarsource2024clean}) inspect source code Abstract Syntax Trees (ASTs) within individual services. They evaluate cyclomatic complexity \cite{mccabe1976complexity}, class cohesion, module coupling (e.g., Lack of Cohesion in Methods [LCOM], Coupling Between Objects [CBO]) \cite{chidamber1994metrics,fenton2014metrics}, and code duplication to flag internal code smells and defect-prone modules \cite{basili1996validation,nagappan2005static,zimmermann2007predicting,menzies2007data}. However, SCA cannot observe runtime communication topology: it does not capture inter-service messaging channels, message broker queue saturation, or cross-host failure propagation.

Recovering system-level structure through static analysis is an active research area. Bushong et al.~\cite{bushong2021static} create communication diagrams from static code analysis, and a recent review looks at nine architecture recovery tools \cite{schneider2024comparison}. In robotics and publish-subscribe systems, tools like HAROS \cite{santos2019haros} use static analysis to find structural problems before launching the system. Together, this research helps build accurate system models for better understanding and drift detection. Unlike these, SaG uses the declared topology to predict how failures might spread.

To close the "Architecture--Code Gap," Static System Analysis (SSA) expands static analysis from individual service code to the whole system architecture. By modeling distributed applications, message topics, brokers, execution nodes, and shared libraries as a connected multigraph, SSA spreads code-level quality metrics across architectural dependencies. As a result, engineering teams can find structural anti-patterns \cite{garcia2009catalogue,taibi2018microservice} and architectural technical debt \cite{li2015technicaldebt} early in the CI/CD process \cite{humble2010continuous,chen2015continuous}, before flawed topologies reach production.

\subsection{Software Quality Models and Multi-Criteria Evaluation}
\label{sec:2.3}

The ISO/IEC 25010:2023 product quality model \cite{iso25010_2023} and the ISO/IEC 25019:2023 Quality-in-Use model \cite{iso25019_2023} define software product quality. From the characteristics listed under Reliability, Maintainability, and Efficiency of Performance in ISO/IEC 25010:2023, SaG focuses on those derived from deployment topology: Availability, Fault Tolerance, Modularity, Modifiability, and Analyzability (Section~\ref{sec:5.1}). Accordingly, topological analysis does not cover other qualities such as Faultlessness, Recoverability, Reusability, and Testability.

Software engineering measures both internal quality, which refers to built-in structural features checked on static artifacts, and external quality, which covers runtime dependability and behavior during system operation \cite{iso25023_2016,iso25021_2012}. In distributed systems, technical debt such as over-centralized message topics or unreplicated brokers can lower internal quality and lead to serious external problems, including performance bottlenecks, queue congestion, and outages.

Combining different structural metrics into a single, auditable quality score is a classic Multi-Criteria Decision Making (MCDM) problem. The Analytic Hierarchy Process (AHP) \cite{saaty1980ahp} uses structured pairwise comparisons and a Consistency Ratio to ensure consistent judgments ($CR \le 0.10$). However, while this ratio finds inconsistencies, it cannot spot matrices filled with preset answers, which this study addresses for the weights (Section~\ref{sec:5.2}). Therefore, AHP provides an audited, explainable Reliability–Maintainability (RM) quality baseline alongside learned graph models.

\subsection{Graph Representation Learning and Explainable AI}
\label{sec:2.4}

Network science uses centrality indices such as degree, closeness, betweenness centrality \cite{freeman1977centrality,brandes2001faster}, articulation points, and PageRank \cite{brin1998pagerank,newman2010networks} to identify important nodes. Key studies on network robustness \cite{albert2000error}, cascading overloads \cite{motter2002cascade}, and interdependent networks \cite{buldyrev2010cascade} model how disruptions spread in connected systems. Percolation models are useful for comparing network breakdowns, but their thresholds assume homogeneous, undirected networks. Because these calculations do not apply to directed multigraphs with different entity types, channels, and QoS contracts, this study uses centrality-based baselines (Section~\ref{sec:6.2}). It suggests multi-type percolation fragmentation as a promising future approach.

Standard network measures have two main problems when applied to software architectures. First, Dimensional Collapse happens when a single centrality score cannot show why a component is critical, such as telling apart a single point of failure, a cascade hub, or an overused library. Second, Semantic Collapse occurs when unweighted metrics treat all nodes and edges the same, mixing up very different architectural elements like message topics, shared libraries, and physical hosts. Together, these problems motivate the core claim that richer architectural analysis is needed.

To overcome the limits of hand-crafted metrics, recent studies use machine learning for network vulnerability analysis and ranking node importance. Examples include FINDER \cite{fan2020keyplayers}, which applies deep reinforcement learning to find key nodes that break network connectivity; DrBC \cite{fan2019learning}, which uses Graph Neural Networks to estimate betweenness centrality in large networks; and PowerGraph \cite{varbella2024powergraph}, which models power grid risks. However, these models assume homogeneous, undirected, unweighted graphs. They cannot be used directly for distributed software systems because Architecture-as-Code manifests are heterogeneous: software topologies are directed multigraphs with several entity types, multiple relation types, and complex QoS contracts. Flattening these architectures into single-layer graphs for DrBC or FINDER causes severe Semantic Collapse and removes important edge contracts. As a result, existing homogeneous models (GCN \cite{kipf2017semi}, GraphSAGE \cite{hamilton2017inductive}, GAT \cite{velickovic2018graph}) average signals across all connection types, which discards relation identity unless it is supplied as a feature. Heterogeneous Graph Neural Networks (RGCN \cite{schlichtkrull2018rgcn}, HAN \cite{wang2019han}, HGT \cite{hu2020hgt}, MAGNN \cite{fu2020magnn}) solve this by using relation-specific transformations. This study uses the Heterogeneous Graph Transformer (HGT) \cite{hu2020hgt} with edge encodings to keep relational semantics when predicting cascade blast radii, and also tests matched homogeneous models (\texttt{GAT-N}, \texttt{GAT-N-QoS}) and closed-form baselines (\texttt{Topo-QoS}) to compare learning approaches (Section~\ref{sec:6.2}). Graph learning has also been used for microservice topologies; for example, Khodabandeh et al.~\cite{khodabandeh2025linkpred} predict future service interactions using graph attention on segmented call graphs. However, that method predicts edge existence based on past interactions. In contrast, this study uses a declared topology as input and predicts the blast radius after removing a node.

A key challenge in using modern AI for software engineering is the black-box problem: deep neural models produce risk scores or embeddings without explaining the underlying causes. In real-world settings, these unclear risk rankings make it hard for developers and site reliability engineers (SREs) to decide whether to replicate hosts, set up circuit breakers, or refactor shared libraries.

Current graph neural network (GNN) explanation methods, such as GNNExplainer \cite{ying2019gnnexplainer} and PGExplainer \cite{luo2020pgexplainer}, identify important subgraphs by masking edges or using parameterized learning. While useful, these methods explain models using internal features instead of standard software engineering terms. SaG addresses this by using a dual-pathway design. Per-relation attention in the HGT pathway is a candidate diagnostic, although on our data its spread across relation types is narrow and largely driven by destination in-degree (Section~\ref{sec:7.3.4}). Meanwhile, the explanation layer links fragility to ISO/IEC quality sub-characteristics (Section~\ref{sec:5}), naming a remediation class for each flagged component.
